% !TEX TS-program = lualatex

%%%
% src::
%   urls = https://tex.stackexchange.com/a/750435/6880 ,
%          https://fr.wikipedia.org/wiki/Icosaèdre ,
%          https://github.com/pfradin/luadraw/discussions/63#discussioncomment-14274453
%%%

\documentclass[border = 3pt]{standalone}

%%%
% Des couleurs '"maison" pour un meilleur design.
%%%
\usepackage{xcolor}

\colorlet{myteal}{teal!50}
\colorlet{mymagenta}{magenta!50}
\colorlet{mylime}{lime!50}

%%%
% La bibliothèque \luadraw allie une facilité d’utilisation
% à un rendu particulièrement soigné.
%%%
\usepackage[3d]{luadraw}

\begin{document}

\begin{luadraw*}{name = icosahedron-by-hand}
local ld = luadraw

local Origin = ld.pt3d.Origin
local M      = ld.pt3d.M

local vecI = ld.pt3d.vecI
local vecJ = ld.pt3d.vecJ
local vecK = ld.pt3d.vecK

------
-- ¨Def de la zone graphique.
------
local graphview = ld.graph3d:new{
  viewdir  = {5, 108},
  size     = {14, 14},
  bbox     = false
}

------
-- Calcul des sommets.
------
local phi = (1 + math.sqrt(5)) / 2

local A1 = M( phi, -1, 0)
local B1 = M( phi,  1, 0)
local C1 = M(-phi,  1, 0)
local D1 = M(-phi, -1, 0)

local A2 = M(0,  phi,  1)
local B2 = M(0,  phi, -1)
local C2 = M(0, -phi, -1)
local D2 = M(0, -phi,  1)

local A3 = M( 1, 0,  phi)
local B3 = M(-1, 0,  phi)
local C3 = M(-1, 0, -phi)
local D3 = M( 1, 0, -phi)

------
-- Construction de l'icosaèdre en trois parties.
--
--
-- tip::
--     Plus bas est donnée une méthode pour ajouter facilement
--     des étiquettes aux sommets, afin de savoir comment définir
--     les différentes faces.
------
local ico_up_facets = {
  {A1, B1, A3},
  {B3, A3, A2},
  {A3, B3, D2},
  {B1, A2, A3},
  {A1, A3, D2}
}

local ico_middle_facets = {
  {B1, A1, D3},
  {C1, D1, B3},
  {B2, A2, B1},
  {A2, B2, C1},
  {D2, C2, A1},
  {C2, D2, D1},
  {A2, C1, B3},
  {D1, D2, B3},
  {B2, B1, D3},
  {A1, C2, D3}
}

local ico_down_facets = {
  {C2, D1, C3},
  {D3, C3, B2},
  {C3, D3, C2},
  {B2, C3, C1},
  {D1, C1, C3}
}

------
-- Pour le tracé des arêtes, nous avons besoin de toutes les faces.
------
local ico_all_facets = ld.facet2poly(
  ld.concat(
    ico_up_facets,
    ico_middle_facets,
    ico_down_facets
  )
)

------
-- Réglage lié aux lignes.
------
ld.Hiddenlinestyle = "dashed"

------
-- Tracé des faces, puis des arêtes.
------
graphview:Dscene3d(
-- Faces colorées (un choix de design).
  graphview:addFacet(
    ico_up_facets,
    {color = "myteal"}
  ),
  graphview:addFacet(
    ico_middle_facets,
    {color = "mymagenta"}
  ),
  graphview:addFacet(
    ico_down_facets,
    {color = "mylime"}
  ),
-- Arêtes cachées tracées en pointillés (un choix de design).
  graphview:addPolyline(
    ld.facetedges(ico_all_facets),
    {
      width  = 12,
      hidden = true
    }
  )
)

------
-- Nous affichons tous les sommets (un choix de design).
--
--
-- note::
--     Pour le moment, l'unique façon de ne pas tracer les sommets
--     cachés est d'effectuer le tracé des points avant celui des
--     faces.
------
graphview:Dballdots3d(
  {
    A1, B1, C1, D1,
    A2, B2, C2, D2,
    A3, B3, C3, D3
  },
  black,
  1.2
)

------
-- Lors du développement, décommenter la ligne suivante ajoute
-- automatiquement des étiquettes aux sommets et aux faces.
------
-- graphview:Dpolynames(ico_all_facets)

------
-- Montrons le résultat de notre oeuvre.
------
graphview:Show()
\end{luadraw*}

\end{document}
